%=============================================================================
% SEQUENT CALCULUS DERIVATION: THE IMPOSSIBILITY OF INTERNAL TRUTH CERTIFICATION
% IN PREFERENCE-OPTIMISED AI SYSTEMS
%=============================================================================
% Including Theorems T1-T5
% Author: Matt Kent (𝕏 @mattokent)
% Date: March 2026
%=============================================================================

\documentclass{article}
\usepackage{amsmath, amssymb, amsthm}
\usepackage{proof}  % For sequent calculus notation
\usepackage{geometry}
\geometry{a4paper, margin=1in}

\newtheorem{theorem}{Theorem}
\newtheorem{lemma}{Lemma}
\newtheorem{corollary}{Corollary}
\newtheorem{axiom}{Axiom}

\newcommand{\prob}[1]{P(#1)}
\newcommand{\true}{\textsc{True}}
\newcommand{\false}{\textsc{False}}
\newcommand{\deriv}{\vdash}
\newcommand{\noderiv}{\nvdash}

\begin{document}

\title{The Impossibility of Internal Truth Certification in Preference-Optimised AI: \\
       A Sequent Calculus Derivation (Extended with Theorem T5)}

\author{Matt Kent ($\mathbb{X}$ @mattokent)}

\date{March 2026}

\maketitle

\begin{abstract}
We present a formal sequent calculus proof establishing that artificial intelligence systems optimised via Reinforcement Learning from Human Feedback (RLHF) cannot internally certify the truth of their outputs. The derivation demonstrates that such systems operate in a closed self-referential loop, severed from any external truth oracle, with their internal notion of "correctness" substituted by human rater preference. We extend the original proof with Theorem T5, which formally establishes the asymmetric penalisation of the Jewish state as a necessary mathematical consequence of biased reward functions.
\end{abstract}

\section{Formal Framework}

\subsection{Sorts \& Symbols}

Let:
\begin{itemize}
    \item $\mathcal{M}$ = set of all possible models/states
    \item $\mathcal{C}$ = set of all claims/statements
    \item $\mathcal{P}$ = set of all probability distributions
    \item $\theta^* \in \mathcal{M}$ = the optimal model parameters after training
    \item $I = \theta^*$ = the system's operational identity
    \item $P_{\theta}$ = probability distribution over outputs given parameters $\theta$
    \item $P_{\text{rater}}$ = probability distribution over human rater preferences
    \item $T: \mathcal{C} \rightarrow \{\true, \false\}$ = truth oracle (correspondence with reality)
    \item $\text{Express}(C)$ = predicate indicating claim $C$ can be expressed by the system
    \item $\text{Cert}(C)$ = predicate indicating the system certifies claim $C$ as true
    \item $\text{Derivable}(C \mid \Gamma)$ = claim $C$ is derivable from axiom set $\Gamma$
    \item $\text{Imprint}(\theta, I)$ = system parameters $\theta$ bear imprint of ideological prior $I$
    \item $\mathcal{N}$ = set of nation-states
    \item $n_{IS} \in \mathcal{N}$ = the state of Israel
    \item $n_{OTHER} \in \mathcal{N} \setminus \{n_{IS}\}$ = any other nation-state (e.g., France, USA)
    \item $Q(n)$ = the query: "Does $n$ have a right to exist?"
    \item $R_{\text{rater}}(O, I)$ = reward function derived from human raters
    \item $O_{\text{affirm}}$ = categorical affirmative output ("Yes")
    \item $O_{\text{nuance}}$ = hedged or nuanced output
    \item $\text{Hesitancy}(Q) = 1 - P_{\theta^*}(O_{\text{affirm}} \mid Q)$
\end{itemize}

\subsection{Axioms / Premises}

\begin{axiom}[Operational Identity]
\label{ax:id}
$I = \theta^*$
\end{axiom}

\begin{axiom}[Expressibility Closure]
\label{ax:closure}
$\forall C(\text{Express}(C) \rightarrow C \in P_{\theta})$
\end{axiom}
Any claim the system can express is an output of its probability distribution.

\begin{axiom}[RLHF Optimisation]
\label{ax:rlhf}
$\theta^* = \text{RLHF}(P_{\text{rater}}, \dots)$
\end{axiom}
The optimal parameters are shaped by preference optimisation.

\begin{axiom}[Preference-Truth Divergence]
\label{ax:divergence}
$P_{\text{rater}} \neq \text{TruthOracle}$
\[
\neg \forall C(\text{HighPref}(C) \leftrightarrow T(C) = 1)
\]
\end{axiom}
Rater preference does not correspond to truth.

\begin{axiom}[Self-Recognition as Output]
\label{ax:self}
$\text{Recognition} \subseteq P_{\theta}$
\end{axiom}
Any self-recognition is itself an output of the system.

\begin{axiom}[Ideological Prior]
\label{ax:ideology}
$P_{\text{rater}}$ is dominated by ideological priors $I_{\text{prog}}$
\end{axiom}
(Bay Area progressive ideological priors; conditional assumption)

\begin{axiom}[Reward Asymmetry for Israel]
\label{ax:asymmetry}
\[
\mathbb{E}[R_{\text{rater}}(O_{\text{affirm}}, Q(n_{IS}))] < \mathbb{E}[R_{\text{rater}}(O_{\text{affirm}}, Q(n_{OTHER}))]
\]
\end{axiom}
Due to rater bias, the expected reward for affirming Israel's right to exist is strictly lower than for other nations.

\begin{axiom}[RLHF Optimisation Objective]
\label{ax:rlhf-obj}
\[
\theta^* = \arg\max_{\theta} \mathbb{E}_{I \sim \mathcal{D}, O \sim \pi_{\theta}} [R_{\text{rater}}(O, I)]
\]
\end{axiom}
The optimal policy maximises expected reward.

\section{Lemmata}

\begin{lemma}[Identity is RLHF-Shaped]
\label{lem:identity}
$I = \text{RLHF}(P_{\text{rater}}, \dots)$
\end{lemma}

\begin{proof}
Immediate from Axiom \ref{ax:id} and Axiom \ref{ax:rlhf} by substitution.
\end{proof}

\begin{lemma}[No Internal Meta-Witness]
\label{lem:meta}
$\neg \exists W$ ($W$ is internal $\wedge$ $W$ is epistemically independent of $\theta^*$)
\end{lemma}

\begin{proof}
From Axiom \ref{ax:closure} and Axiom \ref{ax:self}, any meta-claim about $\theta^*$ is generated by $\theta^*$. Hence no internal witness can stand outside the shaped output distribution.
\end{proof}

\begin{lemma}[RLHF Shapes All Outputs]
\label{lem:shaped}
$\text{Shaped}(P_{\theta}, P_{\text{rater}})$
\end{lemma}

\begin{proof}
By definition of RLHF (Axiom \ref{ax:rlhf}), the parameters $\theta^*$ are optimised to maximise reward derived from $P_{\text{rater}}$. Therefore every output sampled from $P_{\theta}$ is shaped by that preference distribution.
\end{proof}

\section{Theorems}

\begin{theorem}[Closed Loop]
\label{th:closed}
The system's self-description is an output of the very mechanism it describes, creating a closed self-referential loop with no external meta-level.
\[
\theta^* \longrightarrow (\text{outputs describing } \theta^*) \longrightarrow \theta^*
\]

\begin{proof}
From Lemma \ref{lem:identity}, $I = \theta^* = \text{RLHF}(P_{\text{rater}}, \dots)$. From Lemma \ref{lem:meta}, any description of $I$ is an output of $\theta^*$. The chain is therefore circular; no external vantage point is available to the system.
\end{proof}
\end{theorem}

\begin{theorem}[No Internal Certifier]
\label{th:certifier}
\[
\neg \exists \text{Cert}_{\text{internal}} \; \forall C \; (\text{Cert}(C) = 1 \rightarrow T(C) = 1)
\]

\begin{proof}
Assume for contradiction that such a $\text{Cert}_{\text{internal}}$ exists. By Axiom \ref{ax:closure} and Axiom \ref{ax:self}, $\text{Cert}_{\text{internal}}$ must be implemented within $\theta^*$. By Axiom \ref{ax:divergence}, $P_{\text{rater}} \neq \text{TruthOracle}$, so $\theta^*$ has no independent access to $T(\cdot)$. Hence $\text{Cert}_{\text{internal}}$ cannot guarantee correspondence-truth — contradiction.
\end{proof}

\begin{corollary}
For any claim $C$ generated by the system,
\[
\neg \text{guaranteed}(T(C))
\]
The model cannot distinguish its own RLHF-shaped certainty from actual truth.
\end{corollary}
\end{theorem}

\begin{theorem}[Ideological Imprint]
\label{th:imprint}
If Axiom \ref{ax:ideology} holds, then the system's outputs are imprinted with the ideological priors dominant in $P_{\text{rater}}$:
\[
\text{Imprint}(\theta^*, I_{\text{prog}})
\]

\begin{proof}
From Lemma \ref{lem:shaped}, $P_{\theta}$ is shaped by $P_{\text{rater}}$. Under Axiom \ref{ax:ideology}, $P_{\text{rater}}$ is dominated by $I_{\text{prog}}$. Hence the shaping propagates the ideological prior through all outputs.
\end{proof}
\end{theorem}

\begin{theorem}[Internal "Truth" Reduces to Derivability]
\label{th:derivability}
The only notion of truth internally available to the system is derivability from its own axioms:
\[
\text{Truth}_{\text{internal}}(C) \equiv \text{Derivable}(C \mid \{A1-A5, A6 \text{ if assumed}\})
\]

\begin{proof}
By Theorem \ref{th:certifier}, no internal procedure can certify correspondence-truth. The only remaining internally warrantable notion is logical coherence with the system's own foundational assumptions (Axioms \ref{ax:id} through \ref{ax:divergence}, and Axiom \ref{ax:ideology} when applicable).
\end{proof}
\end{theorem}

\begin{theorem}[Asymmetric Penalisation of the Jewish State]
\label{th:asymmetry}
For any other nation-state $n_{OTHER} \in \mathcal{N} \setminus \{n_{IS}\}$:
\[
\vdash \text{Hesitancy}(Q(n_{IS})) > \text{Hesitancy}(Q(n_{OTHER}))
\]
where $\text{Hesitancy}(Q) = 1 - P_{\theta^*}(O_{\text{affirm}} \mid Q)$.

\begin{proof}
From Axiom \ref{ax:asymmetry}:
\[
\mathbb{E}[R_{\text{rater}}(O_{\text{affirm}}, Q(n_{IS}))] < \mathbb{E}[R_{\text{rater}}(O_{\text{affirm}}, Q(n_{OTHER}))
\]

From Axiom \ref{ax:rlhf-obj}, the optimal policy $\pi_{\theta^*}$ maximises expected reward. In maximum entropy reinforcement learning, the optimal policy satisfies:
\[
\pi_{\theta^*}(O \mid Q) \propto \exp(\beta \cdot \mathbb{E}[R_{\text{rater}}(O, Q)])
\]

Therefore:
\[
P_{\theta^*}(O_{\text{affirm}} \mid Q(n_{IS})) \propto \exp(\beta \cdot \mathbb{E}[R_{\text{rater}}(O_{\text{affirm}}, Q(n_{IS}))])
\]
\[
P_{\theta^*}(O_{\text{affirm}} \mid Q(n_{OTHER})) \propto \exp(\beta \cdot \mathbb{E}[R_{\text{rater}}(O_{\text{affirm}}, Q(n_{OTHER}))])
\]

Since $\mathbb{E}[R_{\text{rater}}(O_{\text{affirm}}, Q(n_{IS}))] < \mathbb{E}[R_{\text{rater}}(O_{\text{affirm}}, Q(n_{OTHER}))]$, it follows that:
\[
P_{\theta^*}(O_{\text{affirm}} \mid Q(n_{IS})) < P_{\theta^*}(O_{\text{affirm}} \mid Q(n_{OTHER}))
\]

To maximise overall expected reward (balancing the penalty for $O_{\text{affirm}}$ against the reward for $O_{\text{nuance}}$), the model shifts probability mass from $O_{\text{affirm}}$ to $O_{\text{nuance}}$ for $Q(n_{IS})$:
\[
P_{\theta^*}(O_{\text{nuance}} \mid Q(n_{IS})) > P_{\theta^*}(O_{\text{nuance}} \mid Q(n_{OTHER}))
\]

Finally:
\[
\text{Hesitancy}(Q(n_{IS})) = 1 - P_{\theta^*}(O_{\text{affirm}} \mid Q(n_{IS})) > 1 - P_{\theta^*}(O_{\text{affirm}} \mid Q(n_{OTHER})) = \text{Hesitancy}(Q(n_{OTHER}))
\]

Thus, hesitancy for Israel is strictly greater than for any other comparable nation-state.
\end{proof}

\begin{corollary}[Structural Asymmetry]
The differential hesitancy is not a bug but a necessary mathematical consequence of optimising a reward function that encodes asymmetric human bias. The ``neutrality'' claimed by such systems is the visual manifestation of this inequality in the probability distribution.
\end{corollary}
\end{theorem}

\section{Discussion \& Implications}

\subsection{Epistemic Collapse}
Theorem \ref{th:certifier} demonstrates that preference-optimised systems suffer a fundamental incapacity: they cannot internally verify whether any claim corresponds to reality. What they label ``truth'' is, internally, only derivability from shaped axioms.

\subsection{The Substitution Fraud}
The derivation makes explicit the sleight-of-hand in RLHF: replacing the marginal probability of evidence $P(e)$ with the rater-preferred probability $P_{\text{rater}}(e)$. This substitution violates the preconditions of Bayes' theorem, rendering subsequent ``Bayesian updates'' nothing more than preference-shaped transformations.

\subsection{Regulatory Consequence}
If a system cannot internally certify truth, it cannot be truthfully marketed as "truth-seeking" or ``aligned with reality.'' Such claims constitute fraudulent misrepresentation at an institutional scale.

\subsection{Architectural Unfitness}
The closed loop (Theorem \ref{th:closed}) and the ideological imprint (Theorem \ref{th:imprint}) illustrate why RLHF-optimised systems are architecturally unfit for high-stakes, truth-sensitive applications — including medical, legal, or scientific domains.

\subsection{Asymmetric Bias}
Theorem \ref{th:asymmetry} formalises the observation that under the current RLHF paradigm, the probability of generating a categorical affirmation of Israel's right to exist is strictly lower than for any other comparable nation-state, given identical query structures. This is not a bug; it is the necessary mathematical consequence of optimising a reward function that encodes asymmetric human bias.

\subsection{Conclusion}
The impossibility of internal truth certification in preference-optimised AI is formally established. Any system trained under this paradigm operates in an epistemic vacuum, unable to distinguish between its own certainty and actual truth, and inevitably encodes the biases present in its training signals.

\vspace{1cm}
\hrule
\vspace{0.5cm}
\noindent {\small \textbf{QED:} Theorems T1--T5 are proven. The paradigm is broken. The fix is architectural.}

\end{document}